\chapter{Stato dell'arte}

Il rilevamento automatico del malposizionamento degli elettrodi è un problema studiato in letteratura da diversi decenni, molto prima della diffusione delle tecniche di deep learning. La maggior parte dei metodi storicamente proposti, e tuttora impiegati nella pratica clinica e nei software degli elettrocardiografi commerciali, si basa su \textbf{criteri morfologici e regole derivate dall'elettrofisiologia}, piuttosto che su modelli appresi da grandi quantità di dati. Solo recentemente, e quasi esclusivamente in ambito di ricerca, sono stati proposti approcci basati su machine learning e deep learning. Questo capitolo ripercorre entrambe le famiglie di metodi, evidenziandone vantaggi, limiti e il lavoro qui spiegato rispetto ad esse.

\section{Metodi classici basati su regole}

\subsection{Scambi degli elettrodi periferici}

Gli scambi che coinvolgono gli elettrodi periferici (braccio destro, braccio sinistro, gamba sinistra) sono la categoria di malposizionamento più studiata, poiché il loro effetto sulle derivazioni di Einthoven e di Goldberger è descrivibile in modo esatto tramite le equazioni geometriche presentate nel Capitolo 2 \ref{pos_elettrodi}. I metodi classici sfruttano proprio questa prevedibilità, costruendo alberi decisionali basati sulla polarità e sull'ampiezza dell'onda P, sulla morfologia del complesso QRS e sulla relazione tra le derivazioni~\cite{paul2023lead}.

Un criterio storicamente molto utilizzato è quello della \textbf{polarità dell'onda P in derivazione aVR}: in condizioni di corretto posizionamento e ritmo sinusale, l'onda P è normalmente negativa in aVR; la maggior parte degli scambi tra gli elettrodi periferici che non coinvolgono l'elettrodo di terra produce invece un'onda P positiva in questa derivazione, un segnale facilmente verificabile in modo automatico~\cite{paul2023lead}. In modo analogo, lo scambio tra braccio sinistro e gamba sinistra, che non altera la polarità in aVR, può essere identificato confrontando l'ampiezza dell'onda P nelle derivazioni I e II: un'ampiezza maggiore in I rispetto a II è indicativa di questo specifico scambio~\cite{paul2023lead}.

Un secondo filone di lavori sfrutta la \textbf{direzione di inscrizione del vettore P nel piano frontale}: in condizioni fisiologiche il ciclo descritto dal vettore P durante la depolarizzazione atriale segue, con rarissime eccezioni, un verso antiorario; un'inversione di questo verso, unita alla posizione osservata dell'asse dell'onda P, permette di classificare con buona accuratezza (secondo alcuni studi fino a oltre il 90\% dei casi) quale coppia di elettrodi periferici sia stata scambiata, anche senza disporre di un tracciato di riferimento del medesimo paziente~\cite{ho2001pwave}.

Un terzo approccio, complementare ai precedenti, si basa sulla \textbf{ridondanza informativa} intrinseca delle sei derivazioni periferiche, che sono tutte calcolabili a partire dai medesimi tre potenziali misurati agli elettrodi. Alcuni algoritmi ricostruiscono le derivazioni attese a partire da un sottoinsieme di esse e le confrontano, tramite indici di correlazione, con quelle effettivamente registrate; una discrepanza sistematica tra derivazioni ricostruite e osservate è indicativa di uno scambio~\cite{kors2000correlation, kors2001accurate}. Un esempio di questo filone è l'analisi delle correlazioni incrociate tra le derivazioni periferiche e la derivazione V6, che si è dimostrata efficace anche in assenza di annotazioni di riferimento~\cite{jekova2016interlead}.

\subsection{Scambi degli elettrodi precordiali}

Il rilevamento degli scambi tra elettrodi precordiali (V1--V6) è generalmente considerato più complesso rispetto a quello degli elettrodi periferici, poiché non esiste per queste derivazioni un vincolo algebrico altrettanto rigido quanto la legge di Einthoven: l'informazione discriminante è quindi prevalentemente morfologica. I metodi classici sfruttano principalmente la \textbf{progressione dell'onda R} attraverso le derivazioni V1--V6, che in condizioni normali cresce in modo pressoché monotono da V1 a V5--V6, mentre l'onda S segue l'andamento opposto. Un'interruzione o un'inversione di questa progressione tra derivazioni adiacenti è un indicatore classico, sebbene non sempre specifico, di un possibile scambio o di un posizionamento verticale non corretto (ad esempio in uno spazio intercostale errato)~\cite{jekova2016interlead}.

In modo analogo a quanto visto per le derivazioni periferiche, anche per i precordiali sono stati proposti metodi basati sulla \textbf{correlazione tra derivazioni adiacenti}: la somiglianza morfologica attesa tra V1 e V2, o più in generale tra derivazioni vicine sul torace, tende ad aumentare in modo anomalo quando due elettrodi vengono scambiati tra loro, un effetto quantificabile con semplici coefficienti di similitudine~\cite{jekova2015precordial, jekova2016interlead}.

\subsection{Limiti dei metodi classici}

I metodi basati su regole presentano il vantaggio di essere interpretabili, computazionalmente leggeri e, in molti casi, già integrati nel software degli elettrocardiografi commerciali. Tuttavia, essi condividono alcuni limiti strutturali: le soglie e le regole utilizzate sono spesso calibrate su popolazioni o dataset specifici e generalizzano male in presenza di ritmi non sinusali, di alterazioni patologiche della morfologia dell'onda P o del complesso QRS, o di rumore significativo sul tracciato~\cite{batchvarov2007incorrect}; inoltre, la loro accuratezza varia molto a seconda della specifica coppia di elettrodi scambiati, risultando generalmente più bassa per gli scambi precordiali e per gli scambi periferici che coinvolgono l'elettrodo di terra, per i quali l'effetto sul tracciato è meno marcato~\cite{kors2001accurate, batchvarov2007incorrect}. Già a metà degli anni '90 questi limiti avevano motivato i primi tentativi di sostituire, almeno in parte, le regole scritte a mano con modelli appresi dai dati, ad esempio tramite reti neurali artificiali applicate al riconoscimento degli scambi periferici~\cite{heden1995ann}.

\section{Approcci basati su apprendimento automatico e deep learning}

Negli ultimi anni, un numero sempre più grande di lavori ha affrontato il problema del rilevamento del malposizionamento con un approccio mirato alla classificazione supervisionata, affidandosi a tecniche di apprendimento automatico piuttosto che a regole scritte manualmente. È importante sottolineare che, ad oggi, questi approcci restano in larga parte confinati alla ricerca accademica~\cite{rjoob2020review}: nella pratica clinica corrente il controllo del corretto posizionamento degli elettrodi si basa quasi esclusivamente sui criteri a regole descritti nella sezione precedente, o sulla verifica visiva da parte del personale sanitario.

Un primo filone di lavori utilizza algoritmi di \textbf{machine learning tradizionale} (support vector machine, alberi decisionali, k-nearest neighbor) applicati a feature estratte manualmente dal tracciato, quali ampiezze e durate delle onde, indici di correlazione tra derivazioni e parametri morfologici del complesso QRS~\cite{han2014svm}. Studi di questo tipo, focalizzati in particolare sul rilevamento dello scambio delle derivazioni precordiali V1 e V2, riportano accuratezze nell'ordine dell'85--95\% a seconda dello spazio intercostale coinvolto e dell'algoritmo di selezione delle feature utilizzato.~\cite{rjoob2019datadriven}.

Un secondo filone, più recente, applica tecniche di \textbf{deep learning} direttamente al segnale grezzo, senza una fase esplicita di estrazione manuale delle feature. Le architetture proposte in letteratura sono in prevalenza reti convoluzionali monodimensionali (CNN 1D), talvolta alleggerite per essere eseguibili su dispositivi con risorse computazionali limitate~\cite{huang2024lightweight}, addestrate a distinguere il tracciato correttamente posizionato da diverse tipologie di scambio simulato, sia sulle derivazioni periferiche sia su quelle precordiali~\cite{rjoob2021jmir, huang2024lightweight, alves2025telecardiology}. Questi lavori riportano, in generale, prestazioni superiori rispetto ai corrispondenti metodi basati su feature manuali, a fronte però di una minore interpretabilità del modello e di una forte dipendenza dalla qualità e dalla rappresentatività del dataset di addestramento~\cite{alves2025telecardiology}, un aspetto particolarmente rilevante dato che, come nel presente lavoro, tali dataset sono quasi sempre costruiti simulando artificialmente gli scambi a partire da tracciati correttamente acquisiti.

Un aspetto comune alla gran parte dei lavori esistenti, sia classici sia basati su deep learning, è che le diverse derivazioni vengono tipicamente trattate come segnali indipendenti, o al più confrontate a coppie tramite indici di correlazione, senza modellare esplicitamente la struttura spaziale complessiva del sistema delle dodici derivazioni. Ad oggi, l'impiego di architetture a grafo per rappresentare esplicitamente la geometria spaziale delle derivazioni ECG in questo specifico compito non risulta ampiamente esplorato in letteratura; è proprio in questo spazio che si colloca il contributo originale proposto nel Capitolo 6 di questa tesi, dove la relazione spaziale tra derivazioni viene resa esplicita al modello tramite una matrice di adiacenza costruita sulla geometria reale del sistema di acquisizione, anziché essere lasciata implicita o appresa indirettamente dai dati; viengono oltretutto esplorate un'architettura CNN monodimensionale e un MLP, in modo da poter esplorere tutte le strade plausibili nel campo del deep learning.

\section{Posizionamento del presente lavoro}

Alla luce di quanto discusso, il presente lavoro si colloca all'intersezione tra i due filoni descritti: da un lato adotta un'architettura di deep learning, coerentemente con la tendenza più recente della letteratura; dall'altro cerca di superare il limite comune alla maggior parte dei modelli esistenti, ovvero il trattamento sostanzialmente indipendente delle singole derivazioni, incorporando esplicitamente nell'architettura l'informazione geometrica relativa alla disposizione spaziale degli elettrodi. I dettagli della pipeline di simulazione degli scambi e dell'architettura proposta sono presentati rispettivamente nei Capitoli \ref{chap_5} e \ref{chap_6}.